\section{Introduction}
\label{sec:intro}

\subsection{Widening Gap Between Threat Scale and Defender Capacity}

Modern enterprise environments face an unprecedented deluge of cyber threats that
far outstrips the analytical bandwidth of human defenders. According to official
disclosures by Taiwan's National Security Bureau (NSB) in 2026, external adversaries
launched approximately 2.63 million daily cyber probe and intrusion attempts against
critical national infrastructure in 2025---representing a 6\% year-over-year
escalation~\cite{nsb2026taiwan}. Crucially, these figures quantify externally
originating perimeter scans, port sweeps, and exploit attempts, which must be
rigorously distinguished from internal alert queues processed by corporate Security
Operations Centers (SOCs); confusing external probe volumes with internal SOC alert
loads misrepresents operational queue burdens and may lead to under-resourcing
of frontline analyst teams.

The economic and temporal damages inflicted by uncontained intrusions have reached
catastrophic levels. The IBM Cost of a Data Breach Report 2026 records an all-time
peak global average breach cost of \$4.99 million (+12\% YoY), with an average dwell
and containment lifecycle spanning 247 days~\cite{ibm2026databreach}. Organizations
deploying advanced AI and automation in security workflows reduced breach lifecycles
by 65 days and saved an average of \$1.93 million per incident~\cite{ibm2026databreach}.
Simultaneously, Mandiant M-Trends 2026 reports a global median dwell time of 14 days,
stretching to 122 days for cyber espionage operations~\cite{mandiant2026mtrends}.
These empirical benchmarks demonstrate that detection and containment velocity serves
as the paramount differentiator between minor security incidents and catastrophic
institutional collapse.

\subsection{Alert Fatigue and the Limits of Existing Automation}

A modern SOC relies on Security Information and Event Management (SIEM) and Extended
Detection and Response (XDR) systems to aggregate telemetry across endpoints, network
perimeters, cloud tenants, and identity providers. Tier-1 analysts bear the frontline
responsibility of rapid alert triage: determining whether incoming events represent a
True Positive (TP, confirmed malicious intrusion), a Benign Positive (BP, authorized
administrative activity triggering detection logic), or a False Positive (FP, sensor
erroneous firing)~\cite{alahmadi2022falsepositives}.

The distinction between BP and FP is operationally critical: empirical studies by
Alahmadi et al.~\cite{alahmadi2022falsepositives} establish that the vast majority of
dismissed alarms are benign positives rather than erroneous false alarms. Analysts
map alerts against the MITRE ATT\&CK matrix~\cite{sadlek2025killchain} and adhere to
standard incident handling frameworks such as NIST SP~800-61 Rev.~3~\cite{nist2025sp80061r3},
which mandates prioritized triage based on threat progression rather than chronological
arrival order.

Statistical classifiers such as GBDT-based Copilot Guided Response (CGR)~\cite{freitas2024cgr}
achieve macro-F1 scores above 0.87 on the Microsoft GUIDE benchmark but inherently lack
semantic comprehension of multi-stage attack kill-chains. Conversely, Large Language
Model (LLM) agents provide natural language reasoning capabilities but suffer from a
critical operational failure mode: when an initial log query returns an empty result set,
five evaluated tool-using agent architectures~\cite{tsegai2026aida} prematurely conclude
that no threat exists, yielding a false negative rate (FNR) of at least 40.4\%. The
ExCyTIn-Bench evaluation suite~\cite{wu2026excytin} corroborates this finding, with
investigation rewards plateauing at 0.606 across tested strategies. This gap between
statistical accuracy and semantic investigation reliability constitutes the central
research problem addressed in this work.

\subsection{Research Contributions}

This paper presents four primary contributions to the state of automated SOC alert triage:

\begin{enumerate}
  \item \textbf{Spec-Driven Eight-Module Architecture:} An RFC~2119 behavioral specification
  governing eight decoupled functional components, ensuring auditability and reproducible
  deployment across heterogeneous SOC environments.

  \item \textbf{Bounded ReAct+RAG Control Loop with Evidence Gate:} Algorithm~1
  formalizes a planning loop with an explicit Evidence Gate: when tool queries return
  empty sets, the gate prevents premature FP closure and triggers conditional escalation,
  directly targeting the 40.4\% FNR failure mode.

  \item \textbf{Calibrated Asymmetric Escalation Policy:} A mathematically grounded
  cost-sensitive escalation policy that explicitly models the asymmetric costs between
  missed true-positive attacks and unnecessary human escalations, enabling operators to
  trace a Pareto frontier parameterized by the confidence threshold $\tau_{\text{conf}}$.

  \item \textbf{Dual-Benchmark Reproducible Evaluation Protocol:} A standardized evaluation
  protocol covering both Microsoft GUIDE (707k incidents, 6,100 tenants) and ExCyTIn-Bench
  (57 log tables, 8 attack chains), using GroupKFold partitioning grouped by customer tenant
  ID to prevent cross-institutional data leakage~\cite{arp2022dos}.
\end{enumerate}

The remainder of this paper is organized as follows. Section~II surveys related work in
automated SOC triage, LLM agent evaluation, and human-in-the-loop cybersecurity systems.
Section~III describes the system architecture. Section~IV details the methodology.
Section~V presents the experimental setup. Section~VI reports results. Section~VII
discusses threats to validity, and Section~VIII concludes.
